\documentclass[../../main.tex]{subfiles}
\graphicspath{{../../image/}{../../figure/}}
\begin{document}

\section{证明路线图}\label{section:ch15-证明路线图}

以下我们采用记号 $A(U)=DF(U)$，其特征值记为 $\lambda_i(U)$，右、左特征向量记为 $R_i(U)$、$L_i(U)$，并按 $|R_i(U)|=1$ 与(7.2.3)规范化。

第一步是通过坐标变换 $(x,t)\mapsto(\mu x,\mu t)$ 从\eqref{align:ch15-1513}中消去小参数 $\mu$。事实上，若 $U_\mu$ 是 Cauchy 问题\eqref{align:ch15-1513}、\eqref{align:ch15-1512}的解，则 $U(x,t)=U_\mu(\mu x,\mu t)$ 满足
\begin{align}
\partial_t U(x,t)+A(U(x,t))\partial_xU(x,t)=\partial_x^2U(x,t),\quad-\infty<x<\infty,\ 0<t<\infty,
\label{align:ch15-1521}
\end{align}
其初始条件为
\begin{align}
U(x,0)=U_{0\mu}(x)=U_0(\mu x),\quad-\infty<x<\infty.
\label{align:ch15-1522}
\end{align}

由于 $TV_{(-\infty,\infty)}U_{0\mu}(\cdot)=TV_{(-\infty,\infty)}U_0(\cdot)$，$TV_{(-\infty,\infty)}U_\mu(\cdot,t)=TV_{(-\infty,\infty)}U(\cdot,\mu^{-1}t)$，故只需估计\eqref{align:ch15-1521}的解 $U$ 的总变差即可，其中黏性系数已被缩放为 $1$。关键估计为
\begin{align}
TV_{(-\infty,\infty)}U(\cdot,t)=\|\partial_xU(\cdot,t)\|_{L^1(-\infty,\infty)}<\delta_0,
\label{align:ch15-1523}
\end{align}
其中 $t\in(0,\infty)$，$\delta_0$ 为某个充分小的正数。

上述界来自\eqref{align:ch15-1521}的抛物结构与双曲结构之间的协同作用，方式如下：
\begin{enumerate}
\item\label{item:ch15-152-a} 存在正常数 $\alpha$ 与 $\kappa$，使得当 $TV_{(-\infty,\infty)}U_0(\cdot)<\frac12\kappa\delta_0$ 时，扩散在某个长度为 $\bar t=(\alpha\kappa\delta_0)^{-2}$ 的时间区间 $(0,\bar t]$ 上诱导出\eqref{align:ch15-1523}。进而，当\eqref{align:ch15-1523}在更长的时间区间 $(0,T)$ 上成立、$T>\bar t$ 时，则对任意 $t\in[\bar t,T)$ 有
\begin{align}
\begin{cases}
\|\partial_x^2U(\cdot,t)\|_{L^1(-\infty,\infty)}<2\alpha\delta_0^2,\\
\|\partial_x^3U(\cdot,t)\|_{L^1(-\infty,\infty)}<5\alpha^2\delta_0^3,\\
\|\partial_x^3U(\cdot,t)\|_{L^\infty(-\infty,\infty)}<16\alpha^3\delta_0^4.
\end{cases}
\label{align:ch15-1524}
\end{align}
这将在\cref{section:ch15-扩散效应}中建立。
\item\label{item:ch15-152-b} 当 $t>\bar t$ 时，\eqref{align:ch15-1521}的双曲结构起主导作用，在任一有界或无界、且\eqref{align:ch15-1524}成立的时间区间 $[\bar t,T)$ 上诱导出\eqref{align:ch15-1523}。于是 (b) 与 (a) 联合推出\eqref{align:ch15-1523}对一切 $t\in(0,\infty)$ 成立。
\end{enumerate}

(b) 中结论的验证分若干步进行。\cref{section:ch15-黏性行波分解}说明如何利用\eqref{align:ch15-1521}的 $n$ 个（黏性）行波的叠加，使其最佳拟合区域 $(-\infty,\infty)\times(\bar t,\infty)$ 内任一点 $(x,t)$ 附近解 $U$ 的剖面，从而在一局部坐标系中把 $\partial_xU$ 与 $\partial_tU$ 表示为
\begin{align}
\partial_xU=\sum_{j=1}^nw_jS_j,\qquad\partial_tU=\sum_{j=1}^n\omega_jS_j,
\label{align:ch15-1525}
\end{align}
其分量 $w_j$ 与 $\omega_j$ 满足形如
\begin{align}
\begin{cases}
\partial_tw_j+\partial_x(\sigma_jw_j)-\partial_x^2w_j=\varphi_j,\\
\partial_t\omega_j+\partial_x(\sigma_j\omega_j)-\partial_x^2\omega_j=\psi_j.
\end{cases}
\label{align:ch15-1526}
\end{align}
的标量抛物方程。

下一步在\cref{section:ch15-横波相互作用}、\cref{section:ch15-同族波相互作用}与\cref{section:ch15-能量估计}中进行，即证明：当\eqref{align:ch15-1523}在时间区间 $(0,T)$ 上成立、$T>\bar t$，并且同时
\begin{align}
\int_{\bar t}^{T}\int_{-\infty}^{\infty}\bigl(|\varphi_j(x,t)|+|\phi_j(x,t)|\bigr)\,dx\,dt<\delta_0,\qquad j=1,\cdots,n,
\label{align:ch15-1527}
\end{align}
时，成立更精细的界
\begin{align}
\int_{\bar t}^{T}\int_{-\infty}^{\infty}\bigl(|\varphi_j(x,t)|+|\psi_j(x,t)|\bigr)\,dx\,dt<c\delta_0^2,\qquad j=1,\cdots,n,
\label{align:ch15-1528}
\end{align}
其中 $c$ 与 $T$ 无关。

最后一个要素是标准估计
\begin{align}
\int_{-\infty}^{\infty}|w_j(x,t)|\,dx\leqslant\int_{-\infty}^{\infty}|w_j(x,\bar t)|\,dx+\int_{\bar t}^{t}\int_{-\infty}^{\infty}|\varphi_j(x,\tau)|\,dx\,d\tau,
\label{align:ch15-1529}
\end{align}
它对\eqref{align:ch15-1526}的解在 $t>\bar t$ 时成立。

现可据此证明：当 $\delta_0$ 充分小时，\eqref{align:ch15-1523}对一切 $t\in(0,\infty)$ 成立。设 $TV_{(-\infty,\infty)}U_0(\cdot)<\frac14\kappa\delta_0$。则由上述 (a)，对任意 $t\in(0,\bar t]$ 有 $\|\partial_xU(\cdot,t)\|_{L^1(-\infty,\infty)}<\frac12\delta_0$。现假设\eqref{align:ch15-1523}在某个有界区间 $[\bar t,\bar T]$ 上成立，但在 $t=\bar T$ 处被破坏。对 $c\delta_0<1$，当 $T$ 从 $\bar t$ 增大到 $\bar T$ 时，\eqref{align:ch15-1527}左端除非在更早的某个时刻 $T$ 已取过值 $c\delta_0^2<\delta_0$，否则不可能取值 $\delta_0$。然而这与上述断言“\eqref{align:ch15-1527}蕴含\eqref{align:ch15-1528}”矛盾。因此\eqref{align:ch15-1527}、从而\eqref{align:ch15-1528}对一切 $t\in[\bar t,\bar T]$ 成立。对 $t=\bar T$ 应用\eqref{align:ch15-1529}并利用\eqref{align:ch15-1525}，可得 $\|\partial_xU(\cdot,\bar T)\|_{L^1(-\infty,\infty)}<\frac12\delta_0+c_1\delta_0^2$，当 $2c_1\delta_0<1$ 时它小于 $\delta_0$。于是与“\eqref{align:ch15-1523}在 $t=\bar T$ 处被破坏”的假设矛盾。

稳定性估计\eqref{align:ch15-1515}、\eqref{align:ch15-1516}与\eqref{align:ch15-1517}将在\cref{section:ch15-稳定性估计}中借助\eqref{align:ch15-1523}导出。显然，一旦这些估计建立，便可沿序列 $\{\mu_k\}$（$\mu_k\to0$，$k\to\infty$）取极限，得到\eqref{align:ch15-1511}、\eqref{align:ch15-1512}的解 $U$，并具有稳定性性质\eqref{align:ch15-1518}、\eqref{align:ch15-1519}与\eqref{align:ch15-15110}。随后将证明，任意解 $U=\lim_{k\to\infty}U_{\mu_k}$ 满足 Tame Oscillation Condition 14.10.1。

对具有一般特征族的系统，唯一性问题已在\cref{section:ch15-注记}所引文献中，通过下述程序得到解决。

作为第一步，证明对特殊初始数据 (9.1.12)，

接着证明，\eqref{align:ch15-1511}、\eqref{align:ch15-1512}的任一解 $U=\lim_{k\to\infty}U_{\mu_k}$，在上半平面任一固定点 $(\bar x,\bar t)$ 附近，可由下列二者之一恰当逼近：
\begin{enumerate}
\item\label{item:ch15-152-r-a} Riemann 问题\eqref{align:ch15-1511}的解 $V(x-\bar x,t-\bar t;U_L,U_R)$，其端状态为 $U_L=U(\bar x-,\bar t)$，$U_R=U(\bar x+,\bar t)$；即对任意 $\beta>0$，
\begin{align}
\lim_{h\downarrow0}\frac1h\int_{\bar x-\beta h}^{\bar x+\beta h}|U(x,\bar t+h)-V(x-\bar x,h;U_L,U_R)|\,dx=0;
\label{align:ch15-15210}
\end{align}
\item\label{item:ch15-152-r-b} 线性化系统 Cauchy 问题的解 $W(x-\bar x,t-\bar t)$：
\begin{align}
\begin{cases}
\partial_tW(x,t)+A(U(\bar x,\bar t))\partial_xW(x,t)=0,\\
W(x,0)=U(x-\bar x,\bar t),
\end{cases}
\label{align:ch15-15211}
\end{align}
即存在正常数 $c$ 与 $\beta$，使得对任意 $y<\bar x<z$，
\begin{align}
\limsup_{h\downarrow0}\int_{y+\beta h}^{z-\beta h}|U(x,\bar t+h)-W(x-\bar x,h)|\,dx\leqslant c\bigl[TV_{(y,z)}U(\cdot,\bar t)\bigr]^2.
\label{align:ch15-15212}
\end{align}
\end{enumerate}

可以证明，上述两个条件\eqref{align:ch15-15210}与\eqref{align:ch15-15212}唯一地确定了解，因而整个解族 $\{U_\mu\}$ 当 $\mu\downarrow0$ 时必收敛到 $U$。

\end{document}
